\documentclass[]{article}

\usepackage{amsmath}
\usepackage{multirow}
\usepackage{natbib}
\usepackage{graphicx} 


\begin{document}

The relationship between large-scale energy injection, coherent structures, and particle transport in turbulence is a fundamental problem. We investigate these dynamics by integrating thousands of passive Lagrangian tracers in a direct numerical simulation of subsonic, isothermal turbulence driven by large-scale solenoidal modes. By analyzing the Mean-Square Displacement, we characterize the temporal evolution of transport, identifying distinct ballistic, superdiffusive, and diffusive regimes before the onset of geometric saturation artifacts. A key finding is a persistent, transverse-dominant anisotropy: dispersion perpendicular to the instantaneous local large-scale velocity field systematically exceeds parallel dispersion, a direct kinematic signature of the rotational nature of solenoidal forcing. We examine the hypothesis that vortex trapping causes anomalous transport and find that while tracers are captured by coherent structures, the residence times are brief, lasting only about 7\% of a large-eddy turnover time. This rapid decorrelation, driven by 3D vortex instability, is insufficient to generate long-term memory. Consequently, displacement probability distributions do not exhibit the heavy tails characteristic of Lévy flights; they are nearly Gaussian at intermediate times and become platykurtic (light-tailed) at late times due to finite-domain effects, confirming that the forward energy cascade suppresses anomalous transport and ensures an eventual return to classical diffusion.

\end{document}
                